\documentclass[aspectratio=169,11pt]{beamer}

\usetheme{Madrid}
\usecolortheme{seahorse}
\setbeamertemplate{navigation symbols}{}
\setbeamertemplate{footline}[frame number]

\usepackage[utf8]{inputenc}
\usepackage[T1]{fontenc}
\usepackage[spanish,es-nodecimaldot]{babel}
\usepackage{lmodern}
\usepackage{booktabs}
\usepackage{listings}
\usepackage{xcolor}
\usepackage{hyperref}
\usepackage{ragged2e}

\definecolor{aqlblue}{RGB}{24,90,160}
\definecolor{aqlgreen}{RGB}{0,120,90}
\definecolor{aqlgray}{RGB}{245,245,245}
\definecolor{aqlorange}{RGB}{190,90,20}

\lstdefinelanguage{AQL}{
  morekeywords={FOR,IN,FILTER,RETURN,SORT,LIMIT,LET,COLLECT,WITH,COUNT,INTO,AGGREGATE,INSERT,UPDATE,REPLACE,REMOVE,UPSERT,OPTIONS,OLD,NEW,OUTBOUND,INBOUND,ANY,GRAPH,PRUNE,SHORTEST_PATH,TO,SEARCH,BOOST,ANALYZER,VIEW,DISTINCT,KEEP,AND,OR,NOT,NULL,true,false,ASC,DESC},
  sensitive=true,
  morecomment=[l]{//},
  morecomment=[s]{/*}{*/},
  morestring=[b]",
  morestring=[b]'
}

\lstset{
  language=AQL,
  basicstyle=\ttfamily\scriptsize,
  keywordstyle=\bfseries\color{aqlblue},
  commentstyle=\itshape\color{gray},
  stringstyle=\color{aqlgreen},
  backgroundcolor=\color{aqlgray},
  showstringspaces=false,
  columns=fullflexible,
  keepspaces=true,
  breaklines=true,
  frame=single,
  rulecolor=\color{gray!35},
  tabsize=2
}

\title[Clase AQL]{AQL en ArangoDB: consultas simples, intermedias y avanzadas}
\subtitle{Documentos, agregaciones, subconsultas, grafos y rendimiento}
\author{Alex Di Genova}
\institute{Curso de Bases de Datos / Minería de Datos}
\date{\today}

\begin{document}

%------------------------------------------------
\begin{frame}
  \titlepage
\end{frame}

%------------------------------------------------
\begin{frame}{Objetivos de aprendizaje}
\begin{itemize}
  \item Comprender el rol de AQL como lenguaje declarativo de consulta y manipulación de datos en ArangoDB.
  \item Escribir consultas básicas con \texttt{FOR}, \texttt{FILTER}, \texttt{SORT}, \texttt{LIMIT}, \texttt{LET} y \texttt{RETURN}.
  \item Manipular documentos con \texttt{INSERT}, \texttt{UPDATE}, \texttt{REPLACE}, \texttt{REMOVE} y \texttt{UPSERT}.
  \item Construir consultas intermedias con \texttt{COLLECT}, subconsultas, proyecciones y parámetros.
  \item Aplicar consultas avanzadas sobre grafos, rutas, optimización y buenas prácticas.
\end{itemize}
\end{frame}

%------------------------------------------------
\begin{frame}{¿Qué es AQL?}
\justifying
AQL significa \textbf{ArangoDB Query Language}. Es el lenguaje usado por ArangoDB para consultar y modificar datos almacenados como documentos JSON, colecciones y grafos.

\vspace{0.4cm}
\begin{block}{Idea central}
AQL describe \textbf{qué resultado se quiere obtener}, no necesariamente \textbf{cómo} obtenerlo. Por eso se considera un lenguaje declarativo.
\end{block}

\vspace{0.2cm}
\begin{itemize}
  \item Similar en propósito a SQL, pero con sintaxis adaptada a documentos y grafos.
  \item Es un lenguaje de manipulación de datos: permite leer, insertar, actualizar y eliminar documentos.
  \item No se usa para crear bases, colecciones o índices; esas tareas son administrativas.
\end{itemize}
\end{frame}

%------------------------------------------------
\begin{frame}{Modelo mental: de SQL a AQL}
\begin{table}
\centering
\small
\begin{tabular}{lll}
\toprule
Concepto & SQL & AQL \\
\midrule
Tabla & \texttt{estudiantes} & Colección \texttt{estudiantes} \\
Fila & Registro & Documento JSON \\
Columna & Atributo fijo & Atributo flexible \\
WHERE & \texttt{WHERE} & \texttt{FILTER} \\
SELECT & \texttt{SELECT} & \texttt{RETURN} \\
GROUP BY & \texttt{GROUP BY} & \texttt{COLLECT} \\
JOIN & \texttt{JOIN} & bucles, \texttt{DOCUMENT()}, grafos \\
Relación & FK & documento, edge o referencia \\
\bottomrule
\end{tabular}
\end{table}

\vspace{0.3cm}
\begin{alertblock}{Diferencia clave}
En AQL se trabaja naturalmente con objetos, arreglos y grafos, no solo con tablas planas.
\end{alertblock}
\end{frame}

%------------------------------------------------
\begin{frame}{Caso de estudio}
Usaremos una base de datos pequeña de cursos universitarios:

\vspace{0.2cm}
\begin{itemize}
  \item \texttt{estudiantes}: documentos con nombre, carrera, edad, promedio e intereses.
  \item \texttt{cursos}: documentos con nombre, área y créditos.
  \item \texttt{inscripciones}: edge collection desde estudiante hacia curso, con nota y semestre.
  \item \texttt{sigue}: edge collection estudiante-estudiante, para consultas de red social.
  \item \texttt{requiere}: edge collection curso-curso, para prerrequisitos.
\end{itemize}

\vspace{0.2cm}
\begin{block}{Idea}
El mismo conjunto de datos permite practicar consultas de documentos, agregaciones y grafos.
\end{block}
\end{frame}

%------------------------------------------------
\begin{frame}[fragile]{Preparación: colecciones fuera de AQL}
AQL no crea colecciones. Primero se preparan desde la interfaz web, arangosh o un driver.

\begin{lstlisting}[language=AQL]
// En arangosh, no en AQL puro:
db._create("estudiantes")
db._create("cursos")
db._createEdgeCollection("inscripciones")
db._createEdgeCollection("sigue")
db._createEdgeCollection("requiere")
\end{lstlisting}

\begin{alertblock}{Importante}
Las consultas AQL asumen que las colecciones ya existen.
\end{alertblock}
\end{frame}

%------------------------------------------------
\begin{frame}[fragile]{Datos de ejemplo: estudiantes}
\begin{lstlisting}
FOR doc IN [
  { _key:"e1", nombre:"Ana", edad:21,
    carrera:"Bioinformatica", promedio:6.1,
    activo:true, intereses:["genomica", "R", "grafos"] },
  { _key:"e2", nombre:"Bruno", edad:23,
    carrera:"Ingenieria", promedio:5.4,
    activo:true, intereses:["bases de datos", "Python"] },
  { _key:"e3", nombre:"Carla", edad:20,
    carrera:"Bioinformatica", promedio:6.5,
    activo:true, intereses:["IA", "histologia"] },
  { _key:"e4", nombre:"Diego", edad:24,
    carrera:"Medicina", promedio:4.9,
    activo:false, intereses:["salud", "estadistica"] }
]
INSERT doc INTO estudiantes
\end{lstlisting}
\end{frame}

%------------------------------------------------
\begin{frame}[fragile]{Datos de ejemplo: cursos}
\begin{lstlisting}
FOR doc IN [
  { _key:"bd", nombre:"Bases de Datos", area:"datos", creditos:5 },
  { _key:"prog", nombre:"Programacion", area:"programacion", creditos:6 },
  { _key:"mineria", nombre:"Mineria de Datos", area:"datos", creditos:5 },
  { _key:"bioinfo", nombre:"Bioinformatica", area:"biologia", creditos:4 }
]
INSERT doc INTO cursos
\end{lstlisting}

\vspace{0.2cm}
\begin{block}{Observación}
Cada documento tiene una llave \texttt{\_key}; ArangoDB también genera \texttt{\_id} y \texttt{\_rev}.
\end{block}
\end{frame}

%------------------------------------------------
\begin{frame}[fragile]{Datos de ejemplo: relaciones}
\begin{lstlisting}
FOR doc IN [
  { _from:"estudiantes/e1", _to:"cursos/bd", nota:6.4, semestre:"2026-1" },
  { _from:"estudiantes/e1", _to:"cursos/mineria", nota:6.0, semestre:"2026-1" },
  { _from:"estudiantes/e2", _to:"cursos/bd", nota:5.3, semestre:"2026-1" },
  { _from:"estudiantes/e3", _to:"cursos/mineria", nota:6.7, semestre:"2026-1" },
  { _from:"estudiantes/e4", _to:"cursos/prog", nota:4.5, semestre:"2026-1" }
]
INSERT doc INTO inscripciones

FOR doc IN [
  { _from:"estudiantes/e1", _to:"estudiantes/e2" },
  { _from:"estudiantes/e2", _to:"estudiantes/e3" },
  { _from:"estudiantes/e3", _to:"estudiantes/e4" }
]
INSERT doc INTO sigue
\end{lstlisting}
\end{frame}

%------------------------------------------------
\begin{frame}[fragile]{Datos de ejemplo: prerrequisitos}
\begin{lstlisting}
FOR doc IN [
  { _from:"cursos/mineria", _to:"cursos/bd" },
  { _from:"cursos/mineria", _to:"cursos/prog" },
  { _from:"cursos/bioinfo", _to:"cursos/prog" }
]
INSERT doc INTO requiere
\end{lstlisting}

\vspace{0.25cm}
\begin{block}{Interpretación}
Si existe una arista \texttt{cursos/mineria -> cursos/bd}, entonces Minería de Datos requiere Bases de Datos.
\end{block}
\end{frame}

%------------------------------------------------
\section{Consultas simples}

\begin{frame}[fragile]{1. Primer resultado con RETURN}
Toda consulta de lectura debe producir un resultado con \texttt{RETURN}.

\begin{lstlisting}
RETURN "Hola AQL"
\end{lstlisting}

Resultado esperado:
\begin{lstlisting}
[ "Hola AQL" ]
\end{lstlisting}

\begin{block}{Clave}
El resultado de una consulta AQL siempre es un arreglo, incluso si se devuelve un solo valor.
\end{block}
\end{frame}

%------------------------------------------------
\begin{frame}[fragile]{2. Recorrer una colección con FOR}
\begin{lstlisting}
FOR e IN estudiantes
  RETURN e
\end{lstlisting}

\vspace{0.2cm}
\begin{itemize}
  \item \texttt{FOR e IN estudiantes}: itera por cada documento de la colección.
  \item \texttt{RETURN e}: devuelve el documento completo.
\end{itemize}

\begin{alertblock}{Cuidado}
En colecciones grandes, devolver documentos completos puede ser costoso. Conviene proyectar solo los atributos necesarios.
\end{alertblock}
\end{frame}

%------------------------------------------------
\begin{frame}[fragile]{3. Proyecciones: devolver solo lo necesario}
\begin{lstlisting}
FOR e IN estudiantes
  RETURN {
    nombre: e.nombre,
    carrera: e.carrera,
    promedio: e.promedio
  }
\end{lstlisting}

\vspace{0.2cm}
\begin{block}{Resultado conceptual}
Una lista de objetos más pequeños, no el documento completo.
\end{block}

\begin{lstlisting}
[
  { "nombre":"Ana", "carrera":"Bioinformatica", "promedio":6.1 },
  { "nombre":"Bruno", "carrera":"Ingenieria", "promedio":5.4 }
]
\end{lstlisting}
\end{frame}

%------------------------------------------------
\begin{frame}[fragile]{4. Filtros básicos}
\begin{lstlisting}
FOR e IN estudiantes
  FILTER e.carrera == "Bioinformatica"
  RETURN e.nombre
\end{lstlisting}

\vspace{0.2cm}
Operadores comunes:
\begin{itemize}
  \item Comparación: \texttt{==}, \texttt{!=}, \texttt{<}, \texttt{<=}, \texttt{>}, \texttt{>=}
  \item Pertenencia: \texttt{IN}, \texttt{NOT IN}
  \item Texto: \texttt{LIKE}, \texttt{=\textasciitilde} para expresiones regulares
  \item Lógicos: \texttt{AND}, \texttt{OR}, \texttt{NOT}
\end{itemize}
\end{frame}

%------------------------------------------------
\begin{frame}[fragile]{5. Filtros compuestos}
\begin{lstlisting}
FOR e IN estudiantes
  FILTER e.activo == true
  FILTER e.promedio >= 6.0
  RETURN { nombre: e.nombre, promedio: e.promedio }
\end{lstlisting}

Equivalente:
\begin{lstlisting}
FOR e IN estudiantes
  FILTER e.activo == true AND e.promedio >= 6.0
  RETURN { nombre: e.nombre, promedio: e.promedio }
\end{lstlisting}

\begin{block}{Buena práctica}
Usar varios \texttt{FILTER} puede mejorar la legibilidad. El optimizador decide cómo ejecutar la consulta.
\end{block}
\end{frame}

%------------------------------------------------
\begin{frame}[fragile]{6. Ordenar y limitar}
\begin{lstlisting}
FOR e IN estudiantes
  FILTER e.activo == true
  SORT e.promedio DESC, e.nombre ASC
  LIMIT 3
  RETURN { nombre: e.nombre, promedio: e.promedio }
\end{lstlisting}

\vspace{0.25cm}
\begin{alertblock}{Orden importa}
Usualmente se filtra, luego se ordena y al final se limita. Un \texttt{LIMIT} antes de \texttt{FILTER} puede descartar documentos relevantes.
\end{alertblock}
\end{frame}

%------------------------------------------------
\begin{frame}[fragile]{7. Paginación}
\begin{lstlisting}
FOR e IN estudiantes
  SORT e.nombre ASC
  LIMIT 0, 10
  RETURN e.nombre
\end{lstlisting}

\begin{lstlisting}
FOR e IN estudiantes
  SORT e.nombre ASC
  LIMIT 10, 10
  RETURN e.nombre
\end{lstlisting}

\vspace{0.2cm}
\begin{itemize}
  \item \texttt{LIMIT 0, 10}: primera página.
  \item \texttt{LIMIT 10, 10}: segunda página.
\end{itemize}
\end{frame}

%------------------------------------------------
\begin{frame}[fragile]{8. Variables temporales con LET}
\begin{lstlisting}
FOR e IN estudiantes
  LET estado = e.promedio >= 5.5 ? "aprobado" : "riesgo"
  RETURN {
    nombre: e.nombre,
    promedio: e.promedio,
    estado: estado
  }
\end{lstlisting}

\vspace{0.2cm}
\begin{block}{Uso de LET}
Permite guardar cálculos intermedios, mejorar legibilidad y reutilizar expresiones.
\end{block}
\end{frame}

%------------------------------------------------
\begin{frame}[fragile]{9. Parámetros enlazados}
Los parámetros evitan concatenar texto en las consultas y ayudan a escribir código más seguro.

\begin{lstlisting}
FOR e IN estudiantes
  FILTER e.carrera == @carrera
  FILTER e.promedio >= @min_promedio
  RETURN e
\end{lstlisting}

Parámetros:
\begin{lstlisting}
{
  "carrera": "Bioinformatica",
  "min_promedio": 6.0
}
\end{lstlisting}
\end{frame}

%------------------------------------------------
\begin{frame}[fragile]{10. Parámetro para nombre de colección}
\begin{lstlisting}
FOR doc IN @@coleccion
  LIMIT 5
  RETURN doc
\end{lstlisting}

Parámetros:
\begin{lstlisting}
{
  "@coleccion": "estudiantes"
}
\end{lstlisting}

\begin{alertblock}{Nota}
Para valores se usa \texttt{@param}; para colecciones se usa \texttt{@@param} en la consulta y \texttt{"@param"} en el JSON de parámetros.
\end{alertblock}
\end{frame}

%------------------------------------------------
\section{Consultas intermedias}

\begin{frame}[fragile]{11. Arreglos: expansión}
\begin{lstlisting}
FOR e IN estudiantes
  RETURN {
    nombre: e.nombre,
    intereses: e.intereses,
    primer_interes: e.intereses[0]
  }
\end{lstlisting}

Extraer todos los intereses en una lista plana:
\begin{lstlisting}
RETURN (
  FOR e IN estudiantes
    RETURN e.intereses
)[**]
\end{lstlisting}
\end{frame}

%------------------------------------------------
\begin{frame}[fragile]{12. Filtrar dentro de arreglos}
\begin{lstlisting}
FOR e IN estudiantes
  RETURN {
    nombre: e.nombre,
    intereses_datos: e.intereses[* FILTER
      CONTAINS(LOWER(CURRENT), "datos") OR
      CONTAINS(LOWER(CURRENT), "grafos")
    ]
  }
\end{lstlisting}

\begin{block}{CURRENT}
Dentro de una expresión inline sobre arreglos, \texttt{CURRENT} representa el elemento actual.
\end{block}
\end{frame}

%------------------------------------------------
\begin{frame}[fragile]{13. DOCUMENT(): recuperar por identificador}
\begin{lstlisting}
RETURN DOCUMENT("estudiantes/e1")
\end{lstlisting}

\vspace{0.2cm}
Usando una arista de inscripción:
\begin{lstlisting}
FOR i IN inscripciones
  LET estudiante = DOCUMENT(i._from)
  LET curso = DOCUMENT(i._to)
  RETURN {
    estudiante: estudiante.nombre,
    curso: curso.nombre,
    nota: i.nota
  }
\end{lstlisting}
\end{frame}

%------------------------------------------------
\begin{frame}[fragile]{14. Patrón de join con DOCUMENT()}
Consulta: estudiantes inscritos en Bases de Datos.

\begin{lstlisting}
FOR i IN inscripciones
  FILTER i._to == "cursos/bd"
  LET e = DOCUMENT(i._from)
  RETURN {
    estudiante: e.nombre,
    carrera: e.carrera,
    nota: i.nota
  }
\end{lstlisting}

\begin{block}{Lectura}
La arista tiene \texttt{\_from} y \texttt{\_to}; \texttt{DOCUMENT()} recupera los documentos conectados.
\end{block}
\end{frame}

%------------------------------------------------
\begin{frame}[fragile]{15. Patrón de join con bucles}
\begin{lstlisting}
FOR e IN estudiantes
  FOR i IN inscripciones
    FILTER i._from == e._id
    LET c = DOCUMENT(i._to)
    RETURN {
      estudiante: e.nombre,
      curso: c.nombre,
      nota: i.nota
    }
\end{lstlisting}

\begin{alertblock}{Rendimiento}
Este patrón puede ser costoso si no hay índices adecuados. Para relaciones naturales, una edge collection y traversal pueden ser mejores.
\end{alertblock}
\end{frame}

%------------------------------------------------
\begin{frame}[fragile]{16. Agrupar con COLLECT}
Contar estudiantes por carrera:

\begin{lstlisting}
FOR e IN estudiantes
  COLLECT carrera = e.carrera WITH COUNT INTO n
  SORT n DESC
  RETURN { carrera: carrera, estudiantes: n }
\end{lstlisting}

\vspace{0.25cm}
\begin{block}{Equivalente conceptual}
\texttt{COLLECT} cumple el rol de \texttt{GROUP BY}, pero puede agrupar documentos, calcular agregados y construir arreglos.
\end{block}
\end{frame}

%------------------------------------------------
\begin{frame}[fragile]{17. Agregaciones por grupo}
Promedio de notas por curso:

\begin{lstlisting}
FOR i IN inscripciones
  LET c = DOCUMENT(i._to)
  COLLECT curso = c.nombre
  AGGREGATE
    n = COUNT(),
    promedio = AVERAGE(i.nota),
    maxima = MAX(i.nota)
  SORT promedio DESC
  RETURN {
    curso: curso,
    n: n,
    promedio: ROUND(promedio * 100) / 100,
    maxima: maxima
  }
\end{lstlisting}
\end{frame}

%------------------------------------------------
\begin{frame}[fragile]{18. COLLECT INTO: conservar miembros del grupo}
\begin{lstlisting}
FOR i IN inscripciones
  LET e = DOCUMENT(i._from)
  LET c = DOCUMENT(i._to)
  COLLECT curso = c.nombre INTO grupo = {
    estudiante: e.nombre,
    nota: i.nota
  }
  RETURN {
    curso: curso,
    inscritos: grupo[*].estudiante,
    notas: grupo[*].nota
  }
\end{lstlisting}

\begin{block}{Cuándo usarlo}
Cuando no solo necesitamos estadísticas, sino también los elementos que pertenecen a cada grupo.
\end{block}
\end{frame}

%------------------------------------------------
\begin{frame}[fragile]{19. RETURN DISTINCT vs COLLECT}
\begin{columns}
\column{0.48\textwidth}
Valores únicos simples:
\begin{lstlisting}
FOR e IN estudiantes
  RETURN DISTINCT e.carrera
\end{lstlisting}

\column{0.48\textwidth}
Agrupación flexible:
\begin{lstlisting}
FOR e IN estudiantes
  COLLECT carrera = e.carrera
  RETURN carrera
\end{lstlisting}
\end{columns}

\vspace{0.3cm}
\begin{itemize}
  \item \texttt{RETURN DISTINCT}: simple, para un solo valor.
  \item \texttt{COLLECT}: más flexible; permite contar, agregar y agrupar por varias variables.
\end{itemize}
\end{frame}

%------------------------------------------------
\begin{frame}[fragile]{20. Subconsultas con LET}
Para cada estudiante, listar sus cursos:

\begin{lstlisting}
FOR e IN estudiantes
  LET cursos_tomados = (
    FOR i IN inscripciones
      FILTER i._from == e._id
      LET c = DOCUMENT(i._to)
      SORT c.nombre
      RETURN { curso: c.nombre, nota: i.nota }
  )
  RETURN {
    estudiante: e.nombre,
    n_cursos: LENGTH(cursos_tomados),
    cursos: cursos_tomados
  }
\end{lstlisting}
\end{frame}

%------------------------------------------------
\begin{frame}[fragile]{21. Subconsultas: top-k por estudiante}
Mejores notas de cada estudiante:

\begin{lstlisting}
FOR e IN estudiantes
  LET mejores = (
    FOR i IN inscripciones
      FILTER i._from == e._id
      LET c = DOCUMENT(i._to)
      SORT i.nota DESC
      LIMIT 2
      RETURN { curso: c.nombre, nota: i.nota }
  )
  RETURN { estudiante: e.nombre, mejores: mejores }
\end{lstlisting}

\begin{block}{Idea}
Las subconsultas permiten resolver problemas jerárquicos: un documento principal con resultados anidados.
\end{block}
\end{frame}

%------------------------------------------------
\section{Modificación de datos}

\begin{frame}[fragile]{22. INSERT y RETURN NEW}
\begin{lstlisting}
INSERT {
  _key: "e5",
  nombre: "Elena",
  edad: 22,
  carrera: "Ingenieria",
  promedio: 5.8,
  activo: true
} INTO estudiantes
RETURN NEW
\end{lstlisting}

\begin{block}{NEW}
En consultas de modificación, \texttt{NEW} representa el documento insertado o modificado.
\end{block}
\end{frame}

%------------------------------------------------
\begin{frame}[fragile]{23. UPDATE, REPLACE y REMOVE}
Actualización parcial:
\begin{lstlisting}
UPDATE "e5" WITH { promedio: 6.0 } IN estudiantes
RETURN NEW
\end{lstlisting}

Reemplazo completo:
\begin{lstlisting}
REPLACE "e5" WITH {
  nombre:"Elena", edad:22, carrera:"Ingenieria", activo:true
} IN estudiantes
RETURN NEW
\end{lstlisting}

Eliminar:
\begin{lstlisting}
REMOVE "e5" IN estudiantes
RETURN OLD
\end{lstlisting}
\end{frame}

%------------------------------------------------
\begin{frame}[fragile]{24. UPDATE masivo}
Marcar estudiantes inactivos con bajo promedio:

\begin{lstlisting}
FOR e IN estudiantes
  FILTER e.promedio < 5.0
  UPDATE e WITH { activo: false } IN estudiantes
  RETURN {
    antes: OLD.activo,
    despues: NEW.activo,
    estudiante: NEW.nombre
  }
\end{lstlisting}

\begin{alertblock}{Precaución}
Siempre pruebe primero la parte \texttt{FOR + FILTER + RETURN} antes de ejecutar un \texttt{UPDATE} o \texttt{REMOVE} masivo.
\end{alertblock}
\end{frame}

%------------------------------------------------
\begin{frame}[fragile]{25. UPSERT}
Insertar si no existe; actualizar si ya existe:

\begin{lstlisting}
UPSERT { _key: "e1" }
  INSERT {
    _key: "e1", nombre:"Ana", visitas:1
  }
  UPDATE {
    visitas: OLD.visitas ? OLD.visitas + 1 : 1
  }
IN estudiantes
LET tipo = IS_NULL(OLD) ? "insert" : "update"
RETURN { operacion: tipo, estudiante: NEW.nombre, visitas: NEW.visitas }
\end{lstlisting}

\begin{block}{Uso típico}
Contadores, estados acumulados, sincronización incremental de datos.
\end{block}
\end{frame}

%------------------------------------------------
\section{Consultas avanzadas}

\begin{frame}{Grafos en AQL}
\begin{itemize}
  \item Un grafo tiene \textbf{vértices}: documentos normales.
  \item Y \textbf{aristas}: documentos especiales con \texttt{\_from} y \texttt{\_to}.
  \item AQL permite recorrer grafos con \texttt{OUTBOUND}, \texttt{INBOUND} o \texttt{ANY}.
\end{itemize}

\vspace{0.3cm}
\begin{block}{Sintaxis general}
\texttt{FOR vertice, arista, camino IN min..max DIRECCION inicio edgeCollection RETURN ...}
\end{block}

\vspace{0.2cm}
\begin{itemize}
  \item \texttt{vertice}: nodo visitado.
  \item \texttt{arista}: arista usada para llegar al nodo.
  \item \texttt{camino}: ruta completa con \texttt{vertices} y \texttt{edges}.
\end{itemize}
\end{frame}

%------------------------------------------------
\begin{frame}[fragile]{26. Traversal simple: red social}
Contactos hasta distancia 2 desde Ana:

\begin{lstlisting}
FOR v, e, p IN 1..2 OUTBOUND "estudiantes/e1" sigue
  RETURN {
    contacto: v.nombre,
    profundidad: LENGTH(p.edges),
    camino: p.vertices[*].nombre
  }
\end{lstlisting}

\vspace{0.2cm}
\begin{block}{Interpretación}
\texttt{1..2} significa seguir caminos de mínimo 1 y máximo 2 aristas.
\end{block}
\end{frame}

%------------------------------------------------
\begin{frame}[fragile]{27. Traversal con filtros}
Contactos activos de Ana hasta distancia 2:

\begin{lstlisting}
FOR v, e, p IN 1..2 OUTBOUND "estudiantes/e1" sigue
  FILTER v.activo == true
  SORT LENGTH(p.edges), v.nombre
  RETURN {
    contacto: v.nombre,
    carrera: v.carrera,
    distancia: LENGTH(p.edges)
  }
\end{lstlisting}

\begin{alertblock}{Cuidado}
El filtro se aplica a los vértices emitidos. Si se desea cortar ramas de búsqueda, se usa \texttt{PRUNE}.
\end{alertblock}
\end{frame}

%------------------------------------------------
\begin{frame}[fragile]{28. PRUNE: cortar exploración}
Prerrequisitos de Minería de Datos, cortando cuando se llega a Programación:

\begin{lstlisting}
FOR v, e, p IN 1..3 OUTBOUND "cursos/mineria" requiere
  PRUNE v._key == "prog"
  RETURN {
    prerrequisito: v.nombre,
    camino: p.vertices[*].nombre
  }
\end{lstlisting}

\begin{block}{Diferencia}
\texttt{FILTER} decide qué resultados se devuelven. \texttt{PRUNE} decide qué ramas del grafo se siguen explorando.
\end{block}
\end{frame}

%------------------------------------------------
\begin{frame}[fragile]{29. Opciones de traversal}
\begin{lstlisting}
FOR v, e, p IN 1..3 OUTBOUND "estudiantes/e1" sigue
  OPTIONS {
    order: "bfs",
    uniqueVertices: "path"
  }
  RETURN p.vertices[*].nombre
\end{lstlisting}

\vspace{0.2cm}
\begin{itemize}
  \item \texttt{order: "bfs"}: explora por amplitud.
  \item \texttt{uniqueVertices: "path"}: evita repetir vértices dentro del mismo camino.
  \item Útil para controlar ciclos y el tamaño de la exploración.
\end{itemize}
\end{frame}

%------------------------------------------------
\begin{frame}[fragile]{30. Camino más corto}
\begin{lstlisting}
FOR v, e IN OUTBOUND SHORTEST_PATH
  "estudiantes/e1" TO "estudiantes/e4" sigue
  RETURN v.nombre
\end{lstlisting}

\vspace{0.2cm}
\begin{block}{Pregunta que responde}
¿Cuál es la secuencia mínima de relaciones \texttt{sigue} para llegar desde Ana hasta Diego?
\end{block}

\begin{alertblock}{Nota}
Para grafos con pesos, se pueden usar opciones de peso en variantes de caminos.
\end{alertblock}
\end{frame}

%------------------------------------------------
\begin{frame}[fragile]{31. Recomendación simple usando grafos}
Recomendar contactos de segundo grado que Ana aún no sigue directamente:

\begin{lstlisting}
LET directos = (
  FOR v IN 1..1 OUTBOUND "estudiantes/e1" sigue
    RETURN v._id
)

FOR v, e, p IN 2..2 OUTBOUND "estudiantes/e1" sigue
  FILTER v._id NOT IN directos
  FILTER v._id != "estudiantes/e1"
  RETURN DISTINCT {
    recomendado: v.nombre,
    via: p.vertices[1].nombre
  }
\end{lstlisting}
\end{frame}

%------------------------------------------------
\begin{frame}[fragile]{32. Consulta mixta: documentos + grafos}
Para cada estudiante activo, recuperar cursos y contactos de primer grado:

\begin{lstlisting}
FOR e IN estudiantes
  FILTER e.activo
  LET cursos = (
    FOR i IN inscripciones
      FILTER i._from == e._id
      LET c = DOCUMENT(i._to)
      RETURN c.nombre
  )
  LET contactos = (
    FOR v IN 1..1 OUTBOUND e._id sigue
      RETURN v.nombre
  )
  RETURN { estudiante: e.nombre, cursos: cursos, contactos: contactos }
\end{lstlisting}
\end{frame}

%------------------------------------------------
\begin{frame}[fragile]{33. ArangoSearch: SEARCH en Views}
\texttt{SEARCH} se usa con Views de ArangoSearch, no directamente como \texttt{FILTER} común.

\begin{lstlisting}
FOR doc IN vista_estudiantes
  SEARCH ANALYZER(doc.nombre IN TOKENS(@texto, "text_es"), "text_es")
  SORT BM25(doc) DESC
  LIMIT 10
  RETURN {
    nombre: doc.nombre,
    score: BM25(doc)
  }
\end{lstlisting}

\begin{block}{Cuándo usarlo}
Búsqueda de texto, ranking, scoring, analizadores lingüísticos y consultas tipo buscador.
\end{block}
\end{frame}

%------------------------------------------------
\begin{frame}[fragile]{34. WINDOW: agregados móviles}
Ejemplo conceptual: promedio móvil de notas ordenadas por semestre.

\begin{lstlisting}
FOR i IN inscripciones
  SORT i.semestre ASC
  WINDOW { preceding: 2, following: 0 }
  AGGREGATE promedio_movil = AVERAGE(i.nota)
  RETURN {
    semestre: i.semestre,
    nota: i.nota,
    promedio_movil: promedio_movil
  }
\end{lstlisting}

\begin{block}{Uso típico}
Series temporales, métricas acumuladas, promedios móviles y comparaciones locales.
\end{block}
\end{frame}

%------------------------------------------------
\section{Rendimiento y buenas prácticas}

\begin{frame}[fragile]{35. Revisar planes de ejecución}
En arangosh o desde la interfaz web se puede inspeccionar el plan:

\begin{lstlisting}[language=AQL]
db._explain(`
  FOR e IN estudiantes
    FILTER e.carrera == "Bioinformatica"
    SORT e.promedio DESC
    RETURN e.nombre
`)
\end{lstlisting}

\vspace{0.2cm}
Busque señales como:
\begin{itemize}
  \item \texttt{EnumerateCollectionNode}: posible full scan.
  \item Uso o no uso de índices.
  \item Sorts costosos.
  \item Cantidad estimada de documentos procesados.
\end{itemize}
\end{frame}

%------------------------------------------------
\begin{frame}[fragile]{36. Índices: se crean fuera de AQL}
Ejemplo en arangosh:

\begin{lstlisting}[language=AQL]
db.estudiantes.ensureIndex({
  type: "persistent",
  fields: ["carrera"]
})

// Luego la consulta AQL puede aprovecharlo:
FOR e IN estudiantes
  FILTER e.carrera == "Bioinformatica"
  RETURN e.nombre
\end{lstlisting}

\begin{alertblock}{Regla práctica}
Indexar atributos usados frecuentemente en \texttt{FILTER}, \texttt{SORT} o patrones de búsqueda selectivos.
\end{alertblock}
\end{frame}

%------------------------------------------------
\begin{frame}{Patrones que escalan mejor}
\begin{itemize}
  \item Proyectar solo los atributos necesarios.
  \item Filtrar lo antes posible en la consulta lógica.
  \item Crear índices para filtros frecuentes y de alta selectividad.
  \item Evitar nested loops innecesarios en colecciones grandes.
  \item Usar edge collections y traversals cuando la relación es naturalmente un grafo.
  \item Usar \texttt{COLLECT AGGREGATE} cuando no se necesitan todos los documentos del grupo.
  \item Revisar \texttt{EXPLAIN} y \texttt{PROFILE} antes de optimizar a ciegas.
\end{itemize}
\end{frame}

%------------------------------------------------
\begin{frame}{Errores comunes}
\begin{itemize}
  \item Usar AQL para crear colecciones o índices: eso no corresponde a AQL.
  \item Hacer \texttt{RETURN doc} en colecciones grandes sin necesidad.
  \item Usar \texttt{LIMIT} antes de filtrar u ordenar sin querer.
  \item Olvidar que las subconsultas devuelven arreglos.
  \item Confundir \texttt{FILTER} con \texttt{PRUNE} en traversals.
  \item Ejecutar \texttt{UPDATE} o \texttt{REMOVE} masivos sin probar primero el filtro.
  \item Concatenar parámetros como texto en vez de usar bind parameters.
\end{itemize}
\end{frame}

%------------------------------------------------
\section{Ejercicios}

\begin{frame}{Ejercicios simples}
\begin{enumerate}
  \item Listar los nombres de todos los estudiantes activos.
  \item Listar estudiantes de Bioinformática ordenados por promedio descendente.
  \item Devolver solo nombre, carrera y número de intereses de cada estudiante.
  \item Buscar estudiantes cuyo arreglo \texttt{intereses} contenga \texttt{"grafos"}.
  \item Paginar estudiantes alfabéticamente, 2 por página.
\end{enumerate}
\end{frame}

%------------------------------------------------
\begin{frame}{Ejercicios intermedios}
\begin{enumerate}
  \item Contar estudiantes por carrera con \texttt{COLLECT}.
  \item Calcular promedio de notas por curso.
  \item Para cada estudiante, devolver una lista anidada de cursos inscritos.
  \item Listar cursos con al menos 2 estudiantes inscritos.
  \item Usar \texttt{UPSERT} para registrar una visita por estudiante.
\end{enumerate}
\end{frame}

%------------------------------------------------
\begin{frame}{Ejercicios avanzados}
\begin{enumerate}
  \item Desde \texttt{estudiantes/e1}, listar contactos a distancia 1 y 2.
  \item Recomendar contactos de segundo grado excluyendo contactos directos.
  \item Listar todos los prerrequisitos de \texttt{cursos/mineria} usando traversal.
  \item Encontrar el camino más corto entre dos estudiantes.
  \item Comparar dos planes de ejecución: con y sin índice sobre \texttt{carrera}.
\end{enumerate}
\end{frame}

%%------------------------------------------------
%\begin{frame}[fragile]{Mini evaluación: consulta integradora}
%\textbf{Pregunta:} Para cada carrera, listar estudiantes activos, sus cursos y el promedio de nota por estudiante.
%
%\begin{lstlisting}
%FOR e IN estudiantes
%  FILTER e.activo
%  LET cursos = (
%    FOR i IN inscripciones
%      FILTER i._from == e._id
%      LET c = DOCUMENT(i._to)
%      RETURN { curso: c.nombre, nota: i.nota }
%  )
%  COLLECT carrera = e.carrera INTO grupo = {
%    estudiante: e.nombre,
%    cursos: cursos,
%    promedio_notas: AVERAGE(cursos[*].nota)
%  }
%  RETURN { carrera: carrera, estudiantes: grupo }
%\end{lstlisting}
%\end{frame}

%------------------------------------------------
\begin{frame}{Resumen}
\begin{itemize}
  \item AQL permite consultar documentos y grafos con un mismo lenguaje.
  \item La base de una consulta es \texttt{FOR + FILTER + RETURN}.
  \item \texttt{COLLECT} permite agrupar y agregar datos.
  \item Las subconsultas permiten resultados jerárquicos y consultas más expresivas.
  \item Las edge collections habilitan recorridos de grafos con \texttt{OUTBOUND}, \texttt{INBOUND} y \texttt{ANY}.
  \item La optimización requiere medir: \texttt{EXPLAIN}, \texttt{PROFILE} e índices adecuados.
\end{itemize}
\end{frame}

%------------------------------------------------
\begin{frame}{Referencias principales}
\footnotesize
\begin{itemize}
  \item ArangoDB Documentation. \textit{AQL Documentation}. \url{https://docs.arango.ai/arangodb/stable/aql/}
  \item ArangoDB Documentation. \textit{AQL Data Queries}. \url{https://docs.arango.ai/arangodb/3.12/aql/data-queries/}
  \item ArangoDB Documentation. \textit{High-level AQL operations}. \url{https://docs.arango.ai/arangodb/3.12/aql/high-level-operations/}
  \item ArangoDB Documentation. \textit{COLLECT operation in AQL}. \url{https://docs.arango.ai/arangodb/3.12/aql/high-level-operations/collect/}
  \item ArangoDB Documentation. \textit{Combining queries with subqueries in AQL}. \url{https://docs.arango.ai/arangodb/stable/aql/fundamentals/subqueries/}
  \item ArangoDB Documentation. \textit{Graph queries in AQL}. \url{https://docs.arango.ai/arangodb/stable/aql/graph-queries/}
  \item ArangoDB Documentation. \textit{AQL graph traversals explained}. \url{https://docs.arango.ai/arangodb/stable/aql/graph-queries/traversals-explained/}
\end{itemize}
\end{frame}

\end{document}
